% =============================================================
%  Preámbulo común — TFG Informática
% =============================================================
\usepackage[utf8]{inputenc}
\usepackage[T1]{fontenc}
\usepackage[spanish,es-tabla]{babel}
% IMPORTANTE: cargar amssymb ANTES de newtxmath para evitar el conflicto
% "Command \Bbbk already defined".
\usepackage{amssymb}
\usepackage{newtxtext,newtxmath}

\usepackage[a4paper,top=2cm,bottom=2cm,left=2.5cm,right=2.5cm]{geometry}
\linespread{1.0}

\usepackage{graphicx}
\usepackage{xcolor}
\usepackage{enumitem}
\usepackage{amsmath,mathtools}
\usepackage{booktabs}
\usepackage{caption}
\captionsetup{font=small,labelfont=bf}
\usepackage{listings}
\usepackage{hyperref}
\hypersetup{
  colorlinks=true,
  linkcolor=black,
  citecolor=black,
  urlcolor=blue!50!black
}
\usepackage[backend=biber,style=numeric,sorting=none]{biblatex}
\addbibresource{bib/refs.bib}

\usepackage{fancyhdr}
\pagestyle{fancy}
\fancyhf{}
\fancyhead[L]{\nouppercase{\leftmark}}
\fancyhead[R]{\thepage}
\renewcommand{\headrulewidth}{0.4pt}

\usepackage{titlesec}
\titleformat{\chapter}[hang]{\Large\bfseries}{\thechapter.}{1em}{}
\titlespacing*{\chapter}{0pt}{0pt}{1em}

% Listings: estilo de código.
\definecolor{codebg}{HTML}{F7F7F7}
\definecolor{codekw}{HTML}{0033B3}
\definecolor{codestr}{HTML}{067D17}
\definecolor{codecom}{HTML}{8C8C8C}
\lstdefinestyle{tfgcode}{
  backgroundcolor=\color{codebg},
  basicstyle=\ttfamily\footnotesize,
  keywordstyle=\color{codekw}\bfseries,
  stringstyle=\color{codestr},
  commentstyle=\color{codecom}\itshape,
  breaklines=true,
  showstringspaces=false,
  frame=single,
  framerule=0pt,
  framesep=4pt,
  numbers=left,
  numberstyle=\tiny\color{codecom},
  numbersep=8pt,
  tabsize=2,
  inputencoding=utf8,
  extendedchars=true,
  literate=%
    {á}{{\'a}}1 {é}{{\'e}}1 {í}{{\'i}}1 {ó}{{\'o}}1 {ú}{{\'u}}1
    {Á}{{\'A}}1 {É}{{\'E}}1 {Í}{{\'I}}1 {Ó}{{\'O}}1 {Ú}{{\'U}}1
    {ñ}{{\~n}}1 {Ñ}{{\~N}}1 {¿}{?`}1 {¡}{!`}1 {ü}{{\"u}}1 {Ü}{{\"U}}1,
}
\lstset{style=tfgcode}

% Macros.
\providecommand{\R}{\mathbb{R}}
\providecommand{\E}{\mathbb{E}}
\providecommand{\Prb}{\mathbb{P}}
